\section*{前言}

直觉和直观对于数学学习非常重要，
但这一点在多数线性代数教材中是缺乏的.
在第一章中，
我们暂不急于推出形式化和抽象化的内容，
而是以我们高中就已熟悉的向量概念作为研究对象，
通过观察、总结和猜想，
初步了解线性代数的动机和意义，
建立有效的直觉和直观体系.
在1.1节，
我们引入\textbf{线性映射}的直观概念，
由这一个概念串联起\textbf{矩阵}、\textbf{向量组}和\textbf{线性方程组}
这几个线性代数的重要话题.
从某种程度上说，
这些概念都可以看作是同一回事.
在1.2节，
我们引入\textbf{向量空间}和\textbf{线性子空间}的概念.
这既是我们研究问题的场所，
也是统一矩阵、向量组和线性方程组的工具，
兼具直观性和简洁性.
带着这些工具去学习后面的知识，
会省去很多麻烦.


第一部分（期中考试以前）：
$\mathbb{R}^n$中的线性映射与矩阵